\section{Dataset and Experimental Methodology}
\label{sec:methodology}

\subsection{Task-aligned corpus and families}

The observation unit is a chain/address pair with associated runtime bytecode. The source corpus has four subsets: USENIX-artifact positives (\emph{malicious}); contracts not flagged by the source pipeline (\texttt{benign\_cleared}); general contracts (\texttt{benign\_general}); and a small verified account-abstraction control (\texttt{benign\_AA}). The primary negative set is weak: \texttt{benign\_cleared} means rule-silent, not verified benign, and may contain undetected or dual-use code. The latter two subsets are used only as secondary false-positive controls. Table~\ref{tab:dataset-composition} gives the task-aligned composition.

\begin{table}[t]
  \centering
  \caption{Task-aligned dataset and frozen similarity families at threshold 0.85.}
  \label{tab:dataset-composition}
  \footnotesize
  \setlength{\tabcolsep}{2.3pt}
  \begin{tabular}{@{}p{1.45cm}rrrp{2.25cm}@{}}
    \toprule
    subset & obs. & fam. & sing. & role / label source \\
    \midrule
    malicious & 727 & 209 & 112 & USENIX-artifact positives \\
    \texttt{benign\_cleared} & 1,553 & 635 & 399 & rule-silent weak primary negatives \\
    \texttt{benign\_general} & 797 & 437 & 361 & secondary negatives \\
    \texttt{benign\_AA} & 5 & 5 & 5 & small verified control \\
    \midrule
    total & 3,082 & 1,258 & 856 & retained observations / global families \\
    \bottomrule
  \end{tabular}
  \vspace{2pt}

  \parbox{\columnwidth}{\footnotesize \emph{Note:} ``fam.'' counts families bearing a subset; ``sing.'' counts families with one member of that subset (the total gives global singleton families). A family can bear multiple subsets, so subset family counts must not be summed.}
\end{table}


Before rerunning models, we froze an outcome-blind input policy. Among 76 rows containing 23-byte EIP-7702 designators rather than delegate programs, 32 target runtimes were recovered. Three were retained; 29 were excluded because the same recovered runtime occurred in another frozen family, and 44 unresolved designators were excluded, for 73 designator-source exclusions. We also quarantined all 103 rows belonging to 23 exact-bytecode hashes with conflicting class labels. The policy performs no model-based relabeling and retains no favorable side of a conflict; the resulting corpus has zero cross-class exact-bytecode hashes. Same-class duplicates remain because chain/address is the observation unit: 233 exact groups cover 787 observations.

Related programs are controlled with the pre-existing global family assignment. Bytecode is linearly disassembled into opcode tokens, opcode 4-gram sets receive deterministic 128-permutation BLAKE2b MinHash signatures, and signature agreement estimates Jaccard similarity. Union--find joins pairs at threshold 0.85 across all classes. The task-aligned corpus retains 1,258 frozen families, including 209 with malicious members, 112 with exactly one malicious member, and 856 global singletons. Family IDs and the original outer-fold mapping were preserved after exclusions: the reduced corpus was neither reclustered nor fold-rebalanced. Families are similarity groups, not verified attacker identities. Family-grouped testing controls related-bytecode leakage and provides a more demanding generalization estimate, without proving independence or removing all memorization~\cite{pendlebury2019tesseract}.

\subsection{Evaluation groups and transformations}

We keep four protocol groups separate. \textbf{G-DET} compares malicious with \texttt{benign\_cleared} in five preserved family folds: four folds fit the model and the held-out fold tests it. Each operating threshold maximizes $F_1$ on the fitting-fold predictions. A seeded five-fold random KFold is only a diagnostic, and AUPRC is primary. \textbf{G-MUT} first fits clean M0 models under the same family split, then transforms only held-out positives and reports retained recall over cumulative M0--M3. \textbf{G-VOL} applies the compound metadata/address/selector recipe with variable dead-code flooding to held-out positives; it is a limitation stress test. \textbf{G-ADV} uses outer fold $f$ for test, fold $(f+1)\bmod 5$ for clean-M0 validation, and the other three folds for fitting. Its thresholds maximize validation $F_1$. Values from different groups are not direct before/after comparisons.

The generated inputs follow problem-space evaluation motivation~\cite{pierazzi2020problemspace}, but our guarantee is narrower. M0 is original bytecode. M1 rewrites an existing Solidity-style CBOR metadata region; when none is detected, it appends a synthetic CBOR-like trailer. M2 cumulatively randomizes \texttt{PUSH20} immediates before the original metadata boundary and appends a \texttt{STOP} followed by a deterministic benign-sourced byte chunk sized to approximately 20\% of the original executable-byte length. M3 cumulatively rewrites \texttt{PUSH4} immediates found by the transformed-bytecode scan when they match the sensitive-selector set. These are structure-preserving transformations under our opcode-skeleton checker. For each tier, the checker compares the original and transformed opcode-token sequences only over the original pre-metadata region; all 727 retained positives passed at M1, M2, and M3. This verifies that repository property, not full EVM execution or behavioral equivalence.

F25, F50, F100, and F200 append unreachable benign-sourced byte volumes equal to 25\%, 50\%, 100\%, and 200\% of the original executable-byte length. G-VOL combines metadata, address, and selector transformations with each volume. In G-ADV, by contrast, the F conditions are generated directly from pure M0 and therefore do not include M1--M3. Training sees M0, M1, M2, F25, F50, and F100; M3 and pure-M0 F200 remain held out. Training and test mutations use separate deterministic RNG domains.

\subsection{Augmentation, metrics, and integrity}

G-ADV augments benign and malicious fitting contracts symmetrically. Every source produces the same six training conditions, and each derived sample receives weight $1/n_i$, where $n_i$ is the number of variants for source $i$; all variants from one source therefore contribute one source-contract unit collectively. This source balancing prevents variant count from dominating optimization and blocks simple augmentation-count or padding-label shortcuts, but cannot guarantee the absence of every shortcut. AuthGuard-M0 uses only the original sample from each fitting source; AuthGuard-aug uses the weighted six-condition set. Both retain identical estimator hyperparameters.

AUPRC is primary because the primary corpus is imbalanced; AUROC is secondary. At frozen thresholds we report precision, recall, $F_1$, and false-positive rate. Positive-only G-MUT and G-VOL use retained recall. Main tables aggregate arithmetic means over the five outer test folds. For paired G-ADV differences, uncertainty is estimated with 10,000 replicates that resample frozen test families with replacement, retain every observation in each sampled family, and preserve M0/aug model pairing; these family-clustered intervals are distinct from fold means.

Automated assertions require disjoint source identifiers and families across fit, validation, and test sets; no training/test exact-bytecode-hash overlap; inheritance of each mutant's source family; and no held-out mutation instance in training. All assertions passed in every task-aligned outer fold. They address split and mutation leakage, not label circularity: positives and weak negatives remain scoped to the source artifact and its coverage.
